Se utilizó el simulador {\it CircuitJS} de {\it Falstad} \cite{falstad}
para construir el amperímetro y voltímetro en la Fig. \ref{fig:cositas}.
El simulador no dispone de un galvanómetro, por lo que se utilizó un amperímetro
en su lugar; en una simulación, ambos son idénticos.

Primero, se construyó el galvanómetro real agregando una resistencia $R_g$ en serie,
con un valor de \qty{900}{\ohm} y con una corriente máxima
$I_g^{max}=$\qty{40}{\micro\ampere}. Este circuito se llama {\bf G}. Después, se construyó
un amperímetro de escala máxima $I^{max}=$\qty{4}{\milli\ampere}, 
tomando {\bf G} y agregando la resistencia $R_A$ en paralelo. Este circuito se llama
{\bf A}. Luego, se construyó un voltímetro de escala máxima
$V^{max}=$\qty{12}{\volt}, tomando {\bf G} y agregando la resistencia $R_V$ en
serie. Este circuito se llama {\bf V}.

Posteriormente, para utilizar el amperímetro se construyó un circuito base con una fuente
de voltaje $V=$\qty{5}{\volt} y una resistencia $R=$\qty{2500}{\ohm}. Similarmente, para
utilizar el
voltímetro se construyó un circuito base con una fuente de voltaje $V=$\qty{12}{\volt}
y con dos resistencias $R_1=$\qty{300}{\kilo\ohm} y $R_2=$\qty{100}{\kilo\ohm}, colocando
el voltímetro en los extremos de $R_2$. Después, se cambió la escala máxima de {\bf V} a $V^{max}=$\qty{4}{\volt} en el mismo
circuito y, finalmente, se cambiaron las resistencias $R_1$ y $R_2$ por
\qtylist{30;10}{\kilo\ohm} respectivamente.